\section{Candidate Physical Architectures}
\label{section:synthesize-physical-architectures}

\subsection{Logical Component Alternatives}
Logical components break down the required framework behavior into four building blocks.
Table~\ref{tab:framework-physical-alternatives} gives five alternatives to implement each logical component. For each alternative, a concise description of how it implements each logical component is given.
Valispace Requirements Manager and Computer-Aided Three-Dimensional Interactive Application (CATIA) Stimulus were studied as dedicated MBSE alternatives at the request of the stakeholders, and are considered for all logical components given that they are integrated solutions with ample functionality.
An analysis of Microsoft Excel, Python, the matrix laboratory (MATLAB) environment, and GitHub was also requested for the applicable logical components.
The Concurrent Model-based Engineering Tool (COMET) and the remaining alternatives were chosen from a review of the state of the art in MBSE, consistent with the main objective of this project.

{\small
\setlength{\tabcolsep}{4pt}
\setlength{\parskip}{0pt}
\renewcommand{\arraystretch}{1.02}
\begin{longtable}{@{}>{\raggedright\arraybackslash}p{0.10\linewidth}>{\raggedright\arraybackslash}p{\dimexpr0.27\linewidth-2\tabcolsep\relax}>{\raggedright\arraybackslash}p{\dimexpr0.63\linewidth-2\tabcolsep\relax}@{}}
\caption{Physical implementation alternatives for each logical component with possible implementations of its behavior.}\label{tab:framework-physical-alternatives}\\
\toprule
\textbf{ID} & \textbf{Alternative} & \textbf{Implementation} \\
\midrule
\endfirsthead
\toprule
\textbf{ID} & \textbf{Alternative} & \textbf{Implementation} \\
\midrule
\endhead
\bottomrule
\endfoot
\multicolumn{3}{@{}l@{}}{\textbf{L01 MBSE Workspace}}\\*
L01-A & Valispace Requirements Manager & Represent the mission, CubeSat and subsystems as linked blocks, requirements and computed parameters in Altium Requirements Portal \citep{altiumrequirementsportaldocs}. Tailor properties, modes, interfaces, requirements and verification conditions with custom mappings as needed. Use a custom representational state transfer (REST) or Python API connector to bind L03 analysis inputs and outputs, refresh requirement checks and display numerical results for review \citep{altiumrequirementsportalapi}. Add EPS selection and activity rules alongside Engineering Assistant artificial intelligence (AI) skills \citep{altiumengineeringassistant,altiumrequirementsportalskills}. \\
L01-B & CATIA Stimulus & Use an external editor to define the mission, CubeSat and subsystem hierarchy, and tailor properties, modes, interfaces and requirements. Use a custom connector to bind selected L03 analyses and their inputs and outputs to this model. Use Stimulus to check formalized behavioral requirements and display violations and traces \citep{dassaultstimulus,dassaultstimuluseducation}. Use the external editor to link numerical results to verification conditions for inspection and add EPS component selection and activity guidance. \\
L01-C & Microsoft Excel & Create linked worksheets for the mission, CubeSat, subsystems and EPS components. Tailor properties, modes, interfaces and requirements through structured tables. Use custom macros to bind selected L03 scripts to input and output cells, request analyses and save result links. Use formulas to evaluate verification conditions. Use result sheets to expose measures and failed or unresolved checks for review. Add macros for EPS component selection and activity guidance. \\
L01-D & Python & Use a custom editor to create mission, CubeSat, subsystem and EPS component instances from Python classes. Tailor properties, modes, interfaces and requirements, and bind selected L03 functions to model inputs and result properties. Evaluate verification conditions and display measures, failures and unresolved inputs for review and revision. Add input checks, EPS component selection and activity rules. Consider Nanospace as a model exchange precedent \citep{luccisano2024formalization,review-0340}, although its code is unavailable. \\
L01-E & COMET & Define the mission, CubeSat and subsystem hierarchy in a shared engineering model with linked requirements and numerical parameters \citep{deutsch2022comet}. Tailor properties, modes, interfaces, requirements and verification conditions using custom representations where needed. Use external connectors to bind selected L03 scripts to model inputs and outputs and return numerical results for requirement checks and inspection. Add external EPS selection rules and activity guidance. \\
\midrule
\multicolumn{3}{@{}l@{}}{\textbf{L02 COTS Database}}\\*
L02-A & Valispace Requirements Manager & Populate component hierarchies with catalogue parameters, units and supplier references. Use a custom API adapter to check selected records and supply their identifiers and numerical values to L01 and L03 \citep{altiumrequirementsportaldocs,altiumrequirementsportalapi}. \\
L02-B & CATIA Stimulus & Use an external database for persistent part records, datasheets and identifiers. Use a custom catalogue adapter to resolve L01 selections and supply electrical and dynamic parameter values to mission properties and L03 models. \\
L02-C & Microsoft Excel & Maintain one catalogue row per part with identifiers, units, limits and datasheet links. Use lookup formulas or custom adapters to supply selected, checked parameter values to mission properties and L03 analyses \citep{deutsch2022comet}. \\
L02-D & Python & Provide a custom catalogue with parameter identifiers, units, sources and validation checks. Resolve L01 component references and supply numerical values to mission properties and L03. \\
L02-E & COMET & Use component catalogue entries and datasheet-linked parameters in a shared engineering model \citep{deutsch2022comet}. Use custom mappings to supply selected part properties to L01 and L03. Add external catalogue revision tracking and record validation. \\
\midrule
\multicolumn{3}{@{}l@{}}{\textbf{L03 Simulation and Analysis}}\\*
L03-A & Valispace Requirements Manager & Use selected mission and COTS parameters to evaluate static budgets with linked formulas. Use a custom API connector to run selected dynamic EPS and fault analyses in external solvers and return numerical outputs and execution status to L01 \citep{altiumrequirementsportaldocs,altiumrequirementsportalapi}. \\
L03-B & CATIA Stimulus & Map mission and COTS parameters into external EPS models through custom connectors and the Functional Mock-up Interface \citep{dassaultfmi}. Generate constrained test scenarios and execute selected analyses against formalized requirements. Return numerical outputs, execution status, test traces and requirement coverage to L01 \citep{dassaultstimulus}. \\
L03-C & Microsoft Excel & Use selected mission and COTS inputs in formulas for mass, volume, mode power, eclipse energy and array and battery sizing. Use scenario tables to compare concepts. Use custom adapters to run dynamic electrical and fault analyses externally. Return numerical outputs, result links and execution status to L01 \citep{deutsch2022comet}. \\
L03-D & Python & Run selected custom scripts for budgets, sizing, concept comparisons and electrical and fault calculations using mission and COTS parameters. Return raw numerical outputs and execution status to L01. \\
L03-E & MATLAB/Simulink & Import mission and COTS parameters through a custom connector. Execute selected budgets, sizing and parameter studies in MATLAB, and dynamic EPS and fault models in Simulink. Return numerical outputs and execution status to L01 through the same connector. \\
\midrule
\multicolumn{3}{@{}l@{}}{\textbf{L04 Configuration Management}}\\*
L04-A & Valispace Requirements Manager & Open a working branch in external version control and save a custom mission export for review. Use portal histories and notifications to inspect changes \citep{altiumrequirementsportaldocs,altiumrequirementsportalapi}. Commit the accepted configuration with exact revisions, dependencies, reconstruction instructions, author and authorized decision rationale. \\
L04-B & CATIA Stimulus & Use external version control to branch the mission definition, executable requirements, simulation models and test records \citep{dassaultstimulus}. Save and review changes before committing the accepted baseline with exact revisions, dependencies, reconstruction instructions, author and authorized decision rationale. \\
L04-C & GitHub & Open a Git branch for mission definitions, requirements, catalogue records, models and analyses. Save and inspect the mission changes, then commit the reviewed configuration with author and authorized decision rationale. Use reviewed change requests and tagged releases to save approved baselines, exact revisions, dependencies and reconstruction instructions. \\
L04-D & SeaLion approach & Use Git working branches for mission inputs and generated engineering views, following SeaLion \citep{chiu2023sealion}. Save the mission definition and review its differences before committing the accepted configuration. Extend releases with exact revisions, analysis dependencies, reconstruction instructions, author and authorized decision rationale. \\
L04-E & Nanospace & Extend Nanospace's shared engineering database and dependency links \citep{review-0340} with working branches and saved mission snapshots for review. Commit accepted configurations as immutable baselines with exact revisions, analysis dependencies, reconstruction instructions, author and authorized decision rationale. \\
\end{longtable}
}

\subsection{Synthesized Candidates}

Valispace Requirements Manager (C1) and CATIA Stimulus (C2) were compared as integrated solutions spanning all logical components. Two other candidates were synthesized by combining different alternatives: Version-controlled Python (C3) consists of using native Python for L01-L03, and GitHub for L04, while Version-controlled Excel with Python adapters (C4) uses Microsoft Excel for L01-L02, Python for L03, and GitHub for L04. A fifth candidate representing the as-is capabilities (C5) was introduced to serve as the baseline.

The candidates were compared based on a set of weighted criteria, which have value zero if they match the baseline, +1 to +2 if they improve them, and -1 to -2 if they worsen them. Greater is better.
These criteria support candidate selection based on the stakeholder analysis. All considered candidates are assumed capable of attaining the target MOP values.
The criteria were development complexity (20\%), cost per year (30\%), learning curve (30\%), evidence centralization (10\%) and extensibility (10\%).

\subsection{Candidate Comparison}

Table~\ref{tab:framework-candidate-scores} records each score and its rationale. Lower development complexity, costs and learning effort count as improvements. Greater evidence centralization and extensibility also count as improvements. The weighted total for candidate $j$ is $S_j=\sum_{i=1}^{5}w_i s_{ij}$, where $s_{ij}$ is the criterion score and $w_i$ is its percentage weight expressed as a fraction. C5 has zero for every criterion and a weighted total of zero.

% Weighted totals use the fractional weights, so they also range from -2 to +2.
{\small
\setlength{\tabcolsep}{4pt}
\setlength{\parskip}{0pt}
\renewcommand{\arraystretch}{1.0}
\begin{longtable}{@{}>{\raggedright\arraybackslash}p{0.31\linewidth}>{\centering\arraybackslash}p{\dimexpr0.10\linewidth-2\tabcolsep\relax}>{\centering\arraybackslash}p{\dimexpr0.09\linewidth-2\tabcolsep\relax}>{\raggedright\arraybackslash\rule{0pt}{1.4\baselineskip}}p{\dimexpr0.50\linewidth-2\tabcolsep\relax}@{}}
\caption{Candidate scores and rationales relative to the as-is baseline (C5). Higher is better. C3 is the winning candidate.}\label{tab:framework-candidate-scores}\\
\toprule
\textbf{Criterion} & \textbf{Weight} & \textbf{Score} & \textbf{Rationale} \\
\midrule
\endfirsthead
\toprule
\textbf{Criterion} & \textbf{Weight} & \textbf{Score} & \textbf{Rationale} \\
\midrule
\endhead
\bottomrule
\endfoot
\multicolumn{4}{@{}l@{}}{\textbf{C1 Valispace Requirements Manager}}\\*
Development complexity & 20\% & -2 & Requires custom extensions to cover all requirements, and its interfaces are proprietary and not well documented. \\*
Yearly cost & 30\% & -2 & CA\$1418 yearly. Included in Altium Develop, a professional subscription service. \\*
Learning curve & 30\% & 0 & Includes AI assistance, but the array of controls is large. \\*
Evidence centralization & 10\% & +1 & Depends on the degree of integration allowed by the proprietary interfaces. \\*
Extensibility & 10\% & 0 & Requires a long-term commitment to an expensive product whose usability depends on proprietary interfaces. \\*
\textbf{Weighted total} & \textbf{100\%} & \textbf{-0.90} & \\
\midrule
\multicolumn{4}{@{}l@{}}{\textbf{C2 CATIA Stimulus}}\\*
Development complexity & 20\% & -2 & Requires custom extensions to cover all requirements, and its interfaces are proprietary and not well documented. \\*
Yearly cost & 30\% & -1 & CA\$249 via an individual membership of the International Council on Systems Engineering (INCOSE), for private or academic use only. \\*
Learning curve & 30\% & 0 & Includes AI assistance, but the array of controls is large. \\*
Evidence centralization & 10\% & +1 & Depends on the degree of integration allowed by the proprietary interfaces. \\*
Extensibility & 10\% & 0 & Requires a long-term commitment to an expensive product whose usability depends on proprietary interfaces. 3DEXPERIENCE integration is limited for unauthorized developers. \\*
\textbf{Weighted total} & \textbf{100\%} & \textbf{-0.60} & \\
\midrule
\multicolumn{4}{@{}l@{}}{\textbf{C3 Version-controlled Python}}\\*
Development complexity & 20\% & -1 & A widely used object-oriented programming language. Several publications provide reference patterns, even if code is unavailable. \\*
Yearly cost & 30\% & 0 & No cost to develop, operate, and maintain. \\*
Learning curve & 30\% & 0 & Only the included requirements are developed. Python files can be read and edited by AI agents to provide assistance. Custom assistance can be tailored to user tests. \\*
Evidence centralization & 10\% & +2 & Covers all requirements natively. Includes database management, object-oriented modeling via classes, numerical analysis, and version control with online uploads and authorship traceability. \\*
Extensibility & 10\% & +2 & Ample ecosystem of open-source libraries. \\*
\textbf{Weighted total} & \textbf{100\%} & \textbf{+0.20} & \\
\midrule
\multicolumn{4}{@{}l@{}}{\textbf{C4 Version-controlled Excel with Python adapters}}\\*
Development complexity & 20\% & -1 & A widely used and programmable spreadsheet software. Documentation on its use for systems engineering is scarce. \\*
Yearly cost & 30\% & -1 & CA\$115 yearly for a Microsoft 365 subscription. \\*
Learning curve & 30\% & 0 & Excel is widely used, but macro programming is a specialist topic. \\*
Evidence centralization & 10\% & +1 & Macros require a completely different interface for development, usage, and configuration management.  \\*
Extensibility & 10\% & +1 & The core spreadsheet is proprietary software. The files it generates are difficult to read for version control. \\*
\textbf{Weighted total} & \textbf{100\%} & \textbf{-0.30} & \\
\end{longtable}
}

According to the weighted scores, the winning concept is C3: Version-controlled Python.
The design of this concept is presented in detail in Chapter \ref{ch:design-proposal}, while a detailed economic analysis can be found in Section \ref{section:economic-analysis}.
